% Software scope v0.2 confirmed by the author, 29 Oct 2026.
\ThesisChapterNumbering
\ThesisPart{PART 2. CONTENT}
\chapter{PROBLEM DESCRIPTION}\label{chap:problem}

\section{Detailed Problem Description}
Researchers selecting degron mutations for follow-up experiments face a limited experimental budget.
A 28-residue peptide has 532 possible single-amino-acid substitutions. The software must help the
researcher inspect these candidates, choose a shortlist and retain the evidence needed to revisit
that choice. Its primary deliverable is a traceable prioritisation workflow.

The thesis connects two problems. The research pipeline predicts the effect $\Delta$PSI of a
substitution from a reference peptide, position and replacement amino acid, using experimental
labels and protein-cluster-aware evaluation. The application uses a frozen model to support
analysis, annotation and comparison. The model task remains mutation-level prediction; annotation
and project management do not change its scientific endpoint.

\paragraph{First-release scope.}
The software scope confirmed on 29 October 2026 is one local web application using React with
TypeScript and FastAPI. It includes single-peptide analysis, peptide batches, whole-protein tiling,
ClinVar and approved reference annotations, import of experimental observations for comparison,
CSV/JSON export, and persistent projects and analysis history. The browser, API, local worker and
database belong to the same local workspace. Public hosting, remote deployment and public sharing
are outside acceptance; no separate public and internal editions are required.

The main actor is the researcher operating this local workspace. Account management, multi-user
collaboration, wet-lab/sample management, automatic retraining and clinical diagnosis are outside
scope. Imported measurements remain project data and are not added automatically to training or
benchmark labels. No new biological experiment is required by this thesis.

\section{Functions of the System}
The following use cases describe the target first release, not completed implementation.
Baseline single-tile API/frontend code and a CLI already exist; the added project, job and import
flows still require implementation and verification. The requirements below define the target
behaviour; Chapter~\ref{chap:evaluation} specifies the corresponding acceptance scenarios.

\subsection{Predict the Effect Matrix}
\begin{usecasespec}{Predict the effect matrix}{tab:uc-predict}
  \ucname{Predict effect matrix}{UC-01}
  \ucactor{Researcher}{Important}{Simple}
  \ucfull{Stakeholders and interests:}{Researcher: inspect the effect of every single substitution of a peptide.}
  \ucfull{Summary description:}{Enter a 28-residue WT peptide or select a reference tile; obtain 532 scores in a position-by-amino-acid grid with WT cells masked.}
  \ucfull{Relationships:}{Include: validate input and assess domain. \newline Extend: compare observations, view annotations, save analysis.}
  \ucfull{Event handling:}{1.~Enter/select the peptide. \newline 2.~Validate syntax, coordinates and domain. \newline 3.~Score with a frozen model. \newline 4.~Show domain and warnings before the matrix.}
  \ucfull{Alternative / exceptional flows:}{Invalid input: no scores. Known out-of-domain tile: refuse by default; explicit override records a warning. Unknown, unreliable or borderline domain: score only with a prominent limitation label.}
  \ucfull{Expected outcome:}{A transient preview with raw $\Delta$PSI score, calibrated $P(\Delta\mathrm{PSI}\geq1)$, model identity and domain evidence; explicit save invokes UC-10.}
\end{usecasespec}

\subsection{Rank Substitutions}
\begin{usecasespec}{Rank substitutions}{tab:uc-rank}
  \ucname{Rank substitutions}{UC-02}
  \ucactor{Researcher}{Important}{Simple}
  \ucfull{Stakeholders and interests:}{Researcher: choose a shortlist within an experimental budget.}
  \ucfull{Summary description:}{Sort/filter substitutions within each tile and select candidates.}
  \ucfull{Relationships:}{Include: UC-01 or a completed tile of UC-08.}
  \ucfull{Event handling:}{1.~Open ranking. \newline 2.~Sort by effect/probability and filter by position or replacement. \newline 3.~Select candidates; inspect the top-10\% flag.}
  \ucfull{Alternative / exceptional flows:}{No result: request an analysis first. Empty filter: show an explicit empty state.}
  \ucfull{Expected outcome:}{A tile-specific shortlist retaining coordinates and warnings; no implied cross-protein biological ranking.}
\end{usecasespec}

\subsection{Export Results with Provenance}
\begin{usecasespec}{Export results with provenance}{tab:uc-export}
  \ucname{Export results}{UC-03}
  \ucactor{Researcher}{Important}{Simple}
  \ucfull{Stakeholders and interests:}{Researcher: reuse results in an experimental or computational workflow.}
  \ucfull{Summary description:}{Export a single analysis or complete job as CSV/JSON, including selected annotations/comparisons.}
  \ucfull{Relationships:}{Uses: UC-01, UC-05, UC-06, UC-08 and UC-10.}
  \ucfull{Event handling:}{1.~Choose scope and format. \newline 2.~Write scores, per-item outcomes, coordinates, model/data versions and warnings.}
  \ucfull{Alternative / exceptional flows:}{Export failure preserves the result for retry. Failed/skipped job items remain in the manifest.}
  \ucfull{Expected outcome:}{Complete traceable files, including rows beyond the current page; no public share link.}
\end{usecasespec}

\subsection{Select a Protein and Tiles}
\begin{usecasespec}{Select a protein and tiles}{tab:uc-tile}
  \ucname{Select protein and tiles}{UC-04}
  \ucactor{Researcher}{Important}{Simple}
  \ucfull{Stakeholders and interests:}{Researcher: locate a known protein region without manually copying coordinates.}
  \ucfull{Summary description:}{Search the loaded human UniProt release; choose a tile or analyse the complete protein through tiling.}
  \ucfull{Relationships:}{Continues to UC-01 or UC-08.}
  \ucfull{Event handling:}{1.~Search accession/gene. \newline 2.~Select the reference sequence. \newline 3.~Inspect tiles, domain states and coverage. \newline 4.~Choose the analysis scope.}
  \ucfull{Alternative / exceptional flows:}{No accession, sequence shorter than 28, non-standard residues or unverified mapping: state the reason and affected coverage.}
  \ucfull{Expected outcome:}{Selected reference release and explicit 1-based peptide/protein coordinates.}
\end{usecasespec}

\subsection{Compare with Measured Data}
\begin{usecasespec}{Compare with measured data}{tab:uc-measured}
  \ucname{Compare observations}{UC-05}
  \ucactor{Researcher}{Important}{Simple}
  \ucfull{Stakeholders and interests:}{Researcher: inspect where prediction agrees or disagrees with measurements.}
  \ucfull{Summary description:}{Choose reference observations or a committed project dataset; align exact WT, position and replacement.}
  \ucfull{Relationships:}{Uses: UC-01/08 and UC-09.}
  \ucfull{Event handling:}{1.~Select dataset and replicate/aggregation. \newline 2.~Check mapping, endpoint and units. \newline 3.~Display matched values, exclusions and applicable residuals/metrics.}
  \ucfull{Alternative / exceptional flows:}{Missing/unmatched values are not zero. Different assays remain separate without pooled $\Delta$PSI metrics. Insufficient/constant pairs have explicit undefined-metric states.}
  \ucfull{Expected outcome:}{Comparable prediction/observation views with sample counts and provenance; training observations labelled and imports not automatically called independent validation.}
\end{usecasespec}

\subsection{View ClinVar and Reference Annotations}
\begin{usecasespec}{View ClinVar and reference annotations}{tab:uc-clinvar}
  \ucname{View annotations}{UC-06}
  \ucactor{Researcher}{Important}{Simple}
  \ucfull{Stakeholders and interests:}{Researcher: relate candidate substitutions to published variant and reference evidence.}
  \ucfull{Summary description:}{Show ClinVar missense records, UniProt mapping and approved GPS evidence in release 1.}
  \ucfull{Relationships:}{Extends: UC-01/08; provenance from UC-07.}
  \ucfull{Event handling:}{1.~Enable a layer. \newline 2.~Verify accession/release/WT residue. \newline 3.~Show classification, review/conflict state, source version, attribution and disclaimer.}
  \ucfull{Alternative / exceptional flows:}{Distinguish no match, unmapped custom input, unavailable source and conflicting mapping. Do not guess canonical mappings.}
  \ucfull{Expected outcome:}{Traceable annotations separate from effect predictions; annotation absence does not mean benign.}
\end{usecasespec}

\subsection{View Model and Data Provenance}
\begin{usecasespec}{View model and data provenance}{tab:uc-card}
  \ucname{View model and sources}{UC-07}
  \ucactor{Researcher}{Important}{Simple}
  \ucfull{Stakeholders and interests:}{Researcher and examiner: understand the predicted event, evidence and limitations.}
  \ucfull{Summary description:}{Read frozen model identity, domain, development metrics and source licences/versions.}
  \ucfull{Relationships:}{Supplies evidence to all result/export views.}
  \ucfull{Event handling:}{1.~Open model/data information. \newline 2.~Inspect current metadata or the snapshot associated with a saved run.}
  \ucfull{Alternative / exceptional flows:}{Missing historical artefact: report an integrity problem instead of silently substituting a newer source.}
  \ucfull{Expected outcome:}{Explicit distinction between development evidence, imported comparison and independent validation.}
\end{usecasespec}

\subsection{Analyse Peptide Batches or a Whole Protein}
\begin{usecasespec}{Analyse peptide batches or a whole protein}{tab:uc-batch}
  \ucname{Batch/protein analysis}{UC-08}
  \ucactor{Researcher}{Important}{Complex}
  \ucfull{Stakeholders and interests:}{Researcher: analyse many candidates while retaining item-level failures and context.}
  \ucfull{Summary description:}{Submit bounded peptide records/FASTA or one protein accession/FASTA to a persistent local job.}
  \ucfull{Relationships:}{Uses UC-01/04/10; exports through UC-03.}
  \ucfull{Event handling:}{1.~Choose a project and input mode. \newline 2.~Preview validation, tiling and coverage. \newline 3.~Confirm local persistence and submit. \newline 4.~Track progress and open per-tile results.}
  \ucfull{Alternative / exceptional flows:}{Item failures/refusals preserve successes. Cancellation/restart retain completed results without duplicates. Limits and uncovered residues are explicit.}
  \ucfull{Expected outcome:}{A durable job with reconciled outcome counts, pinned model/source versions and distinct overlapping tile contexts.}
\end{usecasespec}

\subsection{Import Experimental Observations}
\begin{usecasespec}{Import experimental observations}{tab:uc-import}
  \ucname{Import observations}{UC-09}
  \ucactor{Researcher}{Important}{Complex}
  \ucfull{Stakeholders and interests:}{Researcher: compare predictions with locally supplied measurements.}
  \ucfull{Summary description:}{Validate CSV/TSV plus assay, units, source, replicate and reference metadata; commit a dataset revision.}
  \ucfull{Relationships:}{Uses UC-10; supplies UC-05.}
  \ucfull{Event handling:}{1.~Upload into a project. \newline 2.~Preview row-level errors, duplicates, missingness and mapping. \newline 3.~Correct blocking errors and acknowledge warnings. \newline 4.~Explicitly commit.}
  \ucfull{Alternative / exceptional flows:}{Malformed/nonfinite/WT-conflicting records block commit; conflicting duplicates never overwrite. Other endpoints remain separately labelled.}
  \ucfull{Expected outcome:}{An immutable project dataset with file hash and provenance; no update to training data or weights.}
\end{usecasespec}

\subsection{Manage Projects and Analysis History}
\begin{usecasespec}{Manage projects and analysis history}{tab:uc-history}
  \ucname{Manage projects/history}{UC-10}
  \ucactor{Researcher}{Important}{Complex}
  \ucfull{Stakeholders and interests:}{Researcher: revisit the exact analysis used for an earlier experimental decision.}
  \ucfull{Summary description:}{Create/rename/reopen/delete local projects; save and reopen immutable runs and datasets.}
  \ucfull{Relationships:}{Supports UC-01/03/05/08/09.}
  \ucfull{Event handling:}{1.~Create/select project. \newline 2.~Explicitly save a preview or submit/commit work. \newline 3.~Reopen history after restart. \newline 4.~Rerun as a new entry or confirm deletion after a dependency summary.}
  \ucfull{Alternative / exceptional flows:}{Missing snapshot or failed cleanup is explicit; deletion cannot remove shared reference/model data or other projects.}
  \ucfull{Expected outcome:}{Saved input, configuration, results and provenance retained locally until deletion; no account or public-sharing requirement.}
\end{usecasespec}

\section{Non-Functional Requirements}
The detailed identifiers and cases are in the system SRS. The targets below are proposed until measured
on the recorded local computer; no external-host performance or availability result is required.
\subsection{Performance Requirements}
Proposed p95 targets are 300~ms for one tile, 200~ms for listing a 1,000-residue protein's tiles and
1.5~s from click to rendered heatmap. Batch/import resource use is measured separately; bounded jobs,
pagination and cancellation keep the interface usable. Proposed initial bounds are 100 peptides or
one protein of at most 10,000 residues per job, and 10~MiB/100,000 rows per import; these are
configurable design defaults, not measured capacities.
\subsection{Information Safety Requirements}
Serve only approved reference data with source versions and required attribution. ClinVar retains its
review/conflict information and diagnostic disclaimer. Imported observations stay in a project-specific
store with assay/units/replicates and are not automatically pooled into reference or training data.
\subsection{Security Requirements}
Bind application services to loopback and restrict browser origins. Validate all input and uploaded files;
never execute their contents or use supplied names as paths. Unsaved previews are transient. Explicit
saves, job submissions and import commits persist locally until deletion; no sequences or observations
appear in logs, URLs or third-party telemetry. Project history has no automatic 30-day expiry.
\subsection{Reliability Requirements}
Score arrays are deterministic for the same normalised input and frozen configuration. Numerical
agreement with the runner must be below $10^{-9}$ for JSON and $5\times10^{-9}$ for ten-digit CSV.
Saved snapshots survive restart; job recovery does not duplicate completed results. Missing data,
partial failures and invalid mappings remain visible.
\subsection{Maintainability Requirements}
Version models, source snapshots and schema migrations. Document local API/worker/frontend startup and
workspace recovery. Test the clean local workflow rather than cloud installation.
\subsection{Usability and Accessibility Requirements}
Provide keyboard workflows, accessible tables and text/patterns in addition to colour, with WCAG~2.1~AA
as the accessibility target. Desktop use is primary; smaller screens use a scrollable matrix/table fallback.

\section{Implementation and Design Constraints}
\subsection{Implementation Constraints}
One student implements the system within the thesis schedule. React with TypeScript and FastAPI are
fixed choices; PostgreSQL is local project/reference storage and SQLite supports isolated tests/baseline
development. A local worker runs CPU inference. Training and research notebooks remain separate.
Docker may package the same local application; public deployment is not a deliverable.
\subsection{Design Constraints}
Whole-protein analysis retains full 28-residue windows at step 23 and reports uncovered positions.
It does not establish the behaviour of the intact protein. Every result carries assay/domain evidence;
raw $\Delta$PSI, calibrated event probability and clinical annotation are separate quantities. The scope
update does not establish clinical validity, an E3 mechanism or completion of research gates.
